%% sections/02_method.tex

\section{The \AAS\ Integrator}
\label{sec:method}

\subsection{Hamiltonian Framework}
\label{sec:hamiltonian}

We consider the heliocentric Hamiltonian for a test particle,
\begin{equation}
  H(\vq, \vp) = \frac{|\vp|^2}{2}
    - \frac{\mu}{|\vq|}
    + V_{J_2}(\vq)
    + V_{\mathrm{pert}}(\vq, t),
  \label{eq:hamiltonian}
\end{equation}
where $\vq \in \mathbb{R}^3$ is the heliocentric position in AU,
$\vp$ is the conjugate momentum in \unitAUd,
$\mu = k^2 = \mathrm{GMS} = 2.9591\times10^{-4}$\,AU$^3$\,day$^{-2}$
is the Gaussian gravitational constant squared \citep{constants_iau2015},
and $V_{J_2}$ is the solar oblateness correction with
$J_2 = 2.2\times10^{-7}$ \citep{constants_iau2015}.
For the two-body benchmarks in Section~\ref{sec:benchmark},
$V_{J_2} = V_{\mathrm{pert}} = 0$.

\subsection{Yoshida Fourth-Order Composition}
\label{sec:yoshida}

The base map is the Yoshida fourth-order composition
\citep{Yoshida1990} of the leapfrog (Störmer--Verlet) map.
The coefficients are
\begin{equation}
  w_1 = \frac{1}{2 - 2^{1/3}}, \qquad
  w_0 = 1 - 2w_1,
  \label{eq:yoshida_coeffs}
\end{equation}
with $w_1 \approx 1.3512$ and $w_0 \approx -1.7024$.
The composition is the symmetric Drift-Kick-Drift (DKD) sequence,
\begin{equation}
  \Phi_{\Delta s} =
    \begin{aligned}[t]
      &\mathcal{D}(c_1\Delta s)\,
      \mathcal{K}(d_1\Delta s)\,
      \mathcal{D}(c_2\Delta s)\,
      \mathcal{K}(d_2\Delta s)\\
      &\times\,\mathcal{D}(c_2\Delta s)\,
      \mathcal{K}(d_1\Delta s)\,
      \mathcal{D}(c_1\Delta s),
    \end{aligned}
  \label{eq:dkd}
\end{equation}
where $c_1 = w_1/2$, $c_2 = (w_1+w_0)/2$, $d_1 = w_1$,
$d_2 = w_0$. The Drift operator $\mathcal{D}(h)$ advances
$\vq \leftarrow \vq + \vp\,h$; the Kick operator $\mathcal{K}(h)$
advances $\vp \leftarrow \vp + \mathbf{f}(\vq)\,h$, where
$\mathbf{f} = -\nabla\Phi$ is the gravitational force.

\subsection{Hessian-Based Step Size Function}
\label{sec:stepsize}

The Sundman time transformation $dt = g(\vq)\,ds$ maps physical
time $t$ to a pseudo-time $s$, in which the step $\Delta s$ is
fixed while the physical step varies as
\begin{equation}
  \Delta t = g(\vq)\,\Delta s.
  \label{eq:sundman}
\end{equation}

We define the step function as
\begin{equation}
  g(\vq) = \frac{\eps}{\sqrt{\rhoH}},
  \label{eq:g_function}
\end{equation}
where $\eps$ is a dimensionless precision parameter and
$\rhoH \equiv \rho\!\left(\Hess(\vq)\right)$
is the spectral radius of the gravitational potential Hessian
$\Hess = \nabla^2\Phi$.

\subsubsection{Keplerian scaling}
\label{sec:keplerian_scaling}

For the point-mass potential $\Phi = -\mu/r$, the Hessian has
eigenvalues $\{-\mu/r^3,\,-\mu/r^3,\,2\mu/r^3\}$ and
$\rhoH = 2\mu/r^3$. Equation~(\ref{eq:g_function}) then gives
\begin{equation}
  \Delta t = \frac{\eps}{\sqrt{2\mu}}\,r^{3/2}.
  \label{eq:dt_keplerian}
\end{equation}
The $r^{3/2}$ scaling is identical to the optimal adaptive step
derived by \citet{PretoTremaine1999} from a time-transformation
analysis. For a Keplerian orbit the local period scales as
$T_{\mathrm{loc}} \propto r^{3/2}$, so \AAS\ maintains a roughly
constant fraction of the local period per step throughout the orbit.

\subsubsection{J2 correction to the Hessian}
\label{sec:j2_hessian}

When the $J_2$ term is included, the Hessian is augmented as
$\Hess = \Hess_{\mathrm{pm}} + \Hess_{J_2}$, where the
components of $\Hess_{J_2}$ are computed analytically:
\begin{align}
  H_{J_2,ij} &= C\Biggl[
      \left(\frac{15z^2}{r^7} - \frac{3}{r^5}\right)\delta_{ij}
      \notag \\
    &\qquad + q_i\left(\frac{30z\,\delta_{jz}}{r^7}
      - \frac{105z^2\,q_j}{r^9} + \frac{15q_j}{r^7}\right)
      \notag \\
    &\qquad - 6\delta_{iz}\left(\frac{\delta_{jz}}{r^5}
      - \frac{5z\,q_j}{r^7}\right)
    \Biggr],
  \label{eq:j2_hessian}
\end{align}
with $C = \tfrac{1}{2}\mu J_2 R_\odot^2$ and $z = q_3$.

\subsection{Time-Reversibility via Harmonic-Mean Symmetrization}
\label{sec:symmetrization}

To preserve discrete time-reversibility \citep{PretoTremaine1999,
HairerLubichWanner2006}, the step function is symmetrized using
the harmonic mean of $g$ evaluated at the current position
$\vq$ and a predicted position $\vq^* = \vq + \vp\,g(\vq)$:
\begin{equation}
  g_{\mathrm{avg}} =
    \frac{2\,g(\vq)\,g(\vq^*)}{g(\vq) + g(\vq^*)}.
  \label{eq:harmonic_mean}
\end{equation}
One can verify
\begin{equation}
  g_{\mathrm{avg}}(\vq,\vp) = g_{\mathrm{avg}}(\vq^*,-\vp),
\end{equation}
which is the necessary condition
for time-reversal invariance of the discrete map
\citep[\S\,VIII.6]{HairerLubichWanner2006}.

Additional physical bounds are enforced:
(i)~$\Delta t \leq T_{\mathrm{local}}/20$ to avoid step-size
resonances \citep{HernandezZeebeHadden2022}, and
(ii)~$\Delta t \leq 1\,\mathrm{min}$ for $r < 2\,R_\odot$
as a perihelion safeguard.

